% Part: set-theory
% Chapter: ordinals
% Section: introduction

\documentclass[../../../include/open-logic-section]{subfiles}

\begin{document}

\olfileid{sth}{ordinals}{intro}

\olsection{ভূমিকা}

\olref[z][]{chap}-এ আমরা \stagesinf{} নীতির সাহায্যে
ধরে নিয়েছিলাম যে, ক্রমধারার একটি অসীম-তম পর্যায় আছে
(অসীমতার স্বতঃসিদ্ধও দেখো)। কিন্তু \stagessucc{} মেনে
নিলে অসীম-তম পর্যায়েই থামা যায় না; এগিয়ে যেতেই হয়।
তাই প্রথম অসীম পর্যায়ের পরের পর্যায়ে, প্রথম অসীম
পর্যায়ে উপলব্ধ সেটগুলি দিয়ে যত রকম সংগ্রহ গঠন সম্ভব,
সব গঠন করি; তারপর আবার; আবার; আবার; \dots

এখানে অপ্রকাশিতভাবে আমরা সংখ্যার এমন একটি
``স্বজ্ঞাগত'' ধারণা ব্যবহার শুরু করেছি, যার মতে সব
স্বাভাবিক সংখ্যার \emph{পরেও} সংখ্যা থাকতে পারে।
বিশেষত, এখানে ধারণাটি \emph{ট্রান্সফাইনাইট ক্রমসংখ্যা}র।
এই অধ্যায়ের উদ্দেশ্য ধারণাটিকে আরও নির্ভুল করা।
প্রথমে ক্রমসংখ্যার সাধারণ ধারণা আলোচনা করব, তারপর
নির্দিষ্ট কয়েকটি সেটকে আমাদের ক্রমসংখ্যা হিসেবে
স্পষ্টভাবে সংজ্ঞায়িত করব।

\end{document}
